\abstract{

AI agents are increasingly used in mathematics research, but it is often unclear how to use them effectively.
Towards this, we present an extensive case study of how AI was used to improve bounds on the Grothendieck constant $K_G$, which quantifies the hardness between combinatorial problems and their continuous relaxations.
Specifically, while the precise value of $K_G$ is not known, we recently tightened the best known bounds to
\[
    \frac{6\pi}{11}
    \;\le\;
    \KG
    \;\le\;
    \frac{\pi}{2\log(1+\sqrt2)} - 10^{-4}.
\]
Crucially, these improvements were achieved using an AI research system that could arrive at insights deemed novel by domain experts.
We give a detailed discussion of our experience using AI for mathematics research, particularly touching upon its strengths and weaknesses, as well as our experience with creating ideal conditions for AI to arrive at breakthrough insights.
The mathematical results are presented and proved in a companion paper.


% which quantifies the hardness threshold between hard combinatorial problems and their efficient continuous relaxations.


% Points for the final sentence:
% * Documentation of our system
% * Strengths and weaknesses
% * Limitations of current AI
% * Ideal conditions for AI to do discovery
% * "Our harness is just some harness ... but we do wanna talk more generally about harnesses"

% We give an extensive documentation of our system, its strengths and weaknesses, as well as our experinece 



% We give extensive documentation of our system, as well as a discussion on what makes ideal conditions for the AI to discover.




% {
% \color{red}
% Mathematics research is a long-horizon task. It comprises of evolving goals, a growing body of knowledge, and delayed feedback. 
% As AI agents become increasingly common in research, it is important to understand conditions under which agents can successfully collaborate with humans on open ended problems. 
% We report an extended case study of human--AI collaboration on such a task: a six-week, 240-session run of an AI research system that we engineered and steered asynchronously, aimed at the Grothendieck constant ($\KG$). The Grothendieck constant is a central constant in optimization, combinatorial algorithms, and analysis. Despite this,  its value has been a mystery since 1953.
% The collaboration produced novel theoretical machinery that discovered new upper and lower bounds,
% \[
%     \frac{6\pi}{11}
%     \;\le\;
%     \KG
%     \;\le\;
%     \frac{\pi}{2\log(1+\sqrt2)} - 10^{-4},
% \]
% This paper ...
% }


% Mathematics research is a long-horizon task: goals evolve, most attempts fail, and the record of what has been learned outgrows any single model context.
% We report an extended case study of human--AI collaboration on such a task: a six-week, roughly 240-session run of an AI research system that we engineered and steered asynchronously, aimed at the Grothendieck constant $\KG$, a central constant in optimization whose value has been open since 1953.
% The collaboration produced new upper and lower bounds,
% \[
%     \frac{6\pi}{11}
%     \;\le\;
%     \KG
%     \;\le\;
%     \frac{\pi}{2\log(1+\sqrt2)} - 10^{-4},
% \]
% determining the previously unknown tenths digit of $\KG$, together with the new frameworks behind them: asymptotic families of rounding schemes on the upper side, and obstructions valid for all such schemes on the lower side.
% The lower bound was discovered and first proved by the AI system; complete proofs appear in a companion paper.
% This paper explains the mathematics at a high level, documents the methodology, reconstructs the discovery, and analyzes the run's capability profile: strong technical execution, but limited research judgement and fragile maintenance of the research state.
% }

\newpage

\tableofcontents

\newpage
